\chapter{Физички емоции и транзициски договор}
\label{chap:physical-emotions}

\section{Ownership state machine}

Legacy hardware supervisor-от ги именува состојбите прикажани на
\cref{fig:hardware-sm}. Официјалниот continuous runner го задржува истиот
безбедносен принцип: се започнува fail-closed, acquisition се дозволува само
од exact state 1 и по preflight, а fault/release не води кон автоматско
повторно преземање. \code{/expression\_status} е read-only evidence и не
авторизира движење.

\begin{figure}[htbp]
\centering
\resizebox{0.92\textwidth}{!}{\begin{tikzpicture}[node distance=14mm and 18mm]
  \node[box,text width=31mm] (disarmed) {\textbf{DISARMED}\\no SDK owner};
  \node[process,right=of disarmed,text width=31mm] (armed) {\textbf{POSTURE ARMED}\\preflight passed};
  \node[physical,right=of armed,text width=31mm] (pending) {\textbf{ACTION PENDING}\\fresh request};

  \node[physical,below=of pending,text width=31mm] (takeover) {\textbf{SDK TAKEOVER}\\stand once};
  \node[safety,left=of takeover,text width=31mm] (active) {\textbf{ACTION ACTIVE}\\bounded profile};
  \node[process,below=of active,text width=31mm] (recovery) {\textbf{RECOVERY}\\canonical neutral};
  \node[danger,below=of disarmed,text width=31mm] (fault) {\textbf{FAULT}\\hold + release};

  \draw[flow] (disarmed) -- node[above,font=\scriptsize]{state 1 + preflight} (armed);
  \draw[flow] (armed) -- node[above,font=\scriptsize]{authorize} (pending);
  \draw[flow] (pending) -- node[right,font=\scriptsize]{operator start} (takeover);
  \draw[flow] (takeover) -- node[above,font=\scriptsize]{ownership acquired} (active);
  \draw[flow] (active) -- node[right,font=\scriptsize]{change / stale link} (recovery);
  \draw[flow] (recovery.east) to[bend right=28]
    node[right,font=\scriptsize,fill=white,inner sep=1pt]{newest accepted target} (active.south east);
  \draw[flow] (recovery.west) -| node[pos=0.25,below,font=\scriptsize]{settle + release} (disarmed.south);
  \draw[dashedflow] (active.west) --
    node[above,sloped,font=\scriptsize,fill=white,inner sep=1pt]{STOP / hard fault} (fault.east);
  \draw[flow] (fault) -- (disarmed);
\end{tikzpicture}
}
\caption{Логички safety/ownership состојби на одделниот hardware path.}
\label{fig:hardware-sm}
\end{figure}

Normal profile resolver-от е table-driven. Requested EmotionEngine category
се чува непроменета; одделно се објавуваат resolved/active profile и
fallback reason. Така \code{affection} останува видлива како барана емоција,
додека роботот изведува \code{neutral\_animal\_breath}. Прифатена категорија
што не е во runtime allowlist исто добива Neutral, со причина
\code{not\_enabled\_in\_normal\_allowlist}.

\section{Задолжителен exact-neutral transition}

Секој category change го следи \cref{fig:canonical-transition}. Ако paw е
кренат, мора да се спушти и да се потврдат landing и supports. Потоа стариот
loop се запира, командата quintic се враќа во точната canonical stand за 1.5 s,
се држи 0.35 s и се стартува само најновиот pending target. Rapid request го
заменува pending target без да го рестартира return-от; same-category update
може да влијае на следниот loop, но не го прекинува тековниот.

\begin{figure}[htbp]
\centering
\resizebox{\textwidth}{!}{\begin{tikzpicture}[node distance=10mm and 9mm]
  \node[physical,text width=27mm] (old) {category change\\активна фаза};
  \node[safety,right=of old,text width=29mm] (lower) {спушти paw\\потврди landing};
  \node[process,right=of lower,text width=27mm] (cancel) {прекини го\\стариот loop};
  \node[process,right=of cancel,text width=31mm] (return) {quintic return 1.5 s\\exact hold 0.35 s};
  \node[physical,right=of return,text width=28mm] (new) {најнов target\\нова емоција};
  \draw[flow] (old) -- (lower);
  \draw[flow] (lower) -- (cancel);
  \draw[flow] (cancel) -- (return);
  \draw[flow] (return) -- (new);

  \node[danger,below=12mm of cancel,text width=42mm] (hard) {STOP / hard fault\\immediate hold + release};
  \draw[dashedflow] (old.south) |- (hard.west);
  \draw[dashedflow] (lower.south) -- (hard.north west);
  \draw[dashedflow] (new.south) |- (hard.east);
\end{tikzpicture}
}
\caption{Canonical-neutral транзиција и hard-fault bypass.}
\label{fig:canonical-transition}
\end{figure}

STOP и hard faults не чекаат conversational timing. Тие веднаш го активираат
валидираниот hold/release. Stale chat не е unsupported-category fallback: се
извршува safe neutral reset и control се ослободува, наместо robot бесконечно
да дише со застарена состојба.

\section{Neutral}

Neutral е 3.25-секунден three-segment quintic loop: 0.75 s до expanded
endpoint, 1.25 s до compressed endpoint и 1.25 s назад до exact neutral.
HipX и yaw се нула. Base HipY се движи од $-0.008$ rad extension до $+0.020$
rad compression, knee е двојно, roll bias е $\pm0.004$ rad, pitch bias
$\pm0.003$ rad, а combined scalar е $-0.015..+0.027$ rad. Сите четири стапала
остануваат planted.

Физичкиот тест започнал од state 1, battery 37\% и errors 0. По stand и stable
hold биле завршени 27 cycles (87.75 s) до operator handoff, без tracking или
state fault. Максималните spans биле 0.0019074/0.0481415/0.0708771 rad за
HipX/HipY/knee. Четири feedback pauses се опоравиле по 20 frames. Операторот го
прифатил движењето како animal-like breathing.

\section{Joy}

Прифатениот Joy е 5-секунден alternating front-paw gesture: за секое стапало
0.50 s transfer, 0.60 s rise, 0.30 s hold, 0.60 s lower и 0.50 s settle.
Lift target е 50 mm; lateral transfer е 20 mm, а rearward transfer 30 mm лево
и 35 mm десно. Асиметријата е изведена од физички load evidence, не од
естетска претпоставка.

Acceptance suite имал baseline 120.632 N. Front-left unload/landing бил
0.975/36.139 N, а front-right 6.504/30.089 N, со четири supports на двете
landing фази, без safety fault и со clean release. Операторот ја прифатил
анимацијата. Подоцна два normal chat loops и mid-paw retargets поминале.

Природнојазичен Joy turn потоа изложил policy defect. Првиот циклус поминал,
но вториот left unload останал 12.522 N. Paw бил безбедно спуштен на 36.179 N
со четири supports, но runner-от го третирал recoverable miss како terminal и
по release роботот легнал. Корекцијата сега прави 1.5+0.35 s canonical Neutral
и го продолжува Joy под истиот owner само ако landing и четири supports се
потврдени. Сите девет native suites, 29 Python tests, clean build и aarch64
deployment поминале; installed SHA-256 има префикс \code{aa449b}. Не била
испратена motion команда по deployment, па оваа последна корекција е
\status{deployed/offline-tested, live-unvalidated}.

\section{Sadness}

Paw-hover кандидатот не го потврдил unload-от (8.875 N) и бил визуелно
отфрлен. Неколку planted варијанти биле премногу суптилни. Differential bows
од 70 mm и 56 mm стигнале до 10.006 и 10.005 степени pitch и правилно го
активирале hard attitude gate пред heave фазата.

Прифатената варијанта ги спушта само предните corners за 48 mm за 1.50 s,
се стабилизира 0.50 s, изведува три 14 mm heaves со 4 mm alternating side tilt,
држи 1.00 s и се враќа 1.50 s; loop-от е 7.80 s. Единечниот тест и два cycles
за 15.6 s ги задржале четирите supports, немале safety fault и добиле позитивна
операторска оценка. Normal Sadness selection, Sadness$\rightarrow$Neutral и
Sadness$\rightarrow$Fear exact-neutral retarget исто се физички потврдени.

\section{Fear}

Paw-hover Fear candidates не биле прифатени: дел од нив не го минале unload
gate, а операторот оценил дека визуелно не изгледаат како страв. Наместо
bypass, направен е all-feet-planted loop од flinch, rearward recoil, пет
латерални cower transitions, freeze и exact recovery. Три 5-секундни cycles
завршиле за 15 s. Еден stale-feedback pause во stand hold бил задржан и
опоравен; немало safety fault. Финално имало четири supports, state 1/0/0,
battery 60\%, errors 0 и clean release. Операторот потврдил дека Fear работи
правилно. Normal selection, repeat, Fear$\rightarrow$Neutral и
Fear$\rightarrow$Anger retarget се потврдени.

\section{Anger}

Anger не е удар. Тој користи 8 mm brace, 6 mm stance widening, 35 mm paw target,
0.35 s quintic lower, 0.25 s stationary landing dwell и canonical support reset
меѓу paws. Се бара unload, две strong supports, 70 N aggregate load, landing и
четири supports пред следното стапало.

Првиот 15-секунден repeat минал два cycles, но третиот final-right landing
повторно потврдил само три supports. X/Y recenter revision исто не била
доволна. Canonical-reset revision враќа body shift, brace и stance width за 1.0
s, држи 0.35 s, па го проверува истиот gate. Во authorized three-cycle test
сите шест placements поминале; baseline бил 121.912 N, немало safety fault или
feedback pause и release бил чист. Операторот ја прифатил кореографијата.
Normal Anger loops и Anger$\rightarrow$Disgust fallback retarget под еден owner
подоцна поминале.

Како кај Joy, подоцнежен unload miss (6.247 N) бил безбедно спуштен и landed на
28.424 N со четири supports, но старата policy сепак release-ирала. Новата
корекција логира \code{ANGER\_CONTACT\_MISS\_RECOVERED}, прави canonical
Neutral и продолжува Anger само по целосна contact recovery. Девет native
suites, 29 Python tests, clean build и aarch64 deployment поминале; installed
SHA-256 има префикс \code{e72a2d}. Бидејќи немало нов motion test,
корекцијата е \status{deployed/offline-tested, live-unvalidated}.

\section{Неприфатени четири категории}

Affection, Curiosity, Disgust и Surprise имаат специфицирани замисли, но немаат
финални MotionSDK profiles, bounded physical acceptance и операторска оценка.
Нивните нормални барања се зачувуваат како truthful requested emotion, а
resolver-от активира Neutral со
\code{physical\_reaction\_not\_accepted}. Претходни planted prototypes не се
сметаат за прифатени профили. Physical airborne hop, yaw, locomotion и torque
strike остануваат забранети.

Сумарната споредба на деветте категории е дадена во
\cref{tab:physical-emotions}; таа намерно ги одделува прифатените реакции од
планираните профили и од тековниот normal allowlist.

\begin{table}[htbp]
\centering
\caption{Физички статус на деветте категории на 20 септември 2026 година}
\label{tab:physical-emotions}
\small
\begin{tabularx}{\textwidth}{L{0.13\textwidth}YL{0.31\textwidth}}
\toprule
\textbf{Емоција} & \textbf{Физички профил} & \textbf{Точно доказно ниво} \\
\midrule
Neutral & 3.25 s animal-like breathing, сите стапала на подлогата & Физички извршен и операторски прифатен; canonical waypoint. \\
Joy & 5 s alternating 50 mm front-paw gesture & Профилот и chat transitions се физички прифатени; persistent contact-miss recovery е deployed/offline-tested, но live revalidation недостига. \\
Sadness & 48 mm front-only planted bow со три heaves & Физички извршен, repeat и chat retarget поминати, операторски прифатен. \\
Anger & canonical-reset alternating 35 mm controlled paw placements & Профилот, repeats и retarget се физички потврдени; последната persistent miss-recovery е deployed/offline-tested, но live revalidation недостига. \\
Fear & planted flinch/recoil/cower/freeze & Физички извршен, chat selection/retarget поминати, операторски прифатен. \\
Affection & Neutral fallback & Сопствен профил е планиран, неимплементиран и нефизички прифатен. \\
Curiosity & Neutral fallback & Сопствен профил е планиран, неимплементиран и нефизички прифатен. \\
Disgust & Neutral fallback & Сопствен профил е планиран, неимплементиран и нефизички прифатен. \\
Surprise & Neutral fallback & Сопствен профил е планиран, неимплементиран и нефизички прифатен. \\
\bottomrule
\end{tabularx}
\end{table}


И покрај физички прифатените пет реакции, checked-in normal allowlist останува
само \code{neutral}. Причините се pending consolidated operator verdict за
целосниот sweep и live revalidation на последните Joy/Anger persistent-recovery
корекции. Додавање environment value не е доказ за commissioning.
